\chapter*{摘\quad 要}
\addcontentsline{toc}{chapter}{摘要}

本文建立一份与浙江工业大学研究生学位论文格式规范相匹配的硕士论文 \LaTeX{} 示例。当前摘要为版面占位文字，用于验证中文小四号正文、1.25 倍行距、页眉页码和关键词格式。正式写作时，应在此概括研究目的、研究方法、主要结果及结论，硕士论文中文摘要建议控制在学校要求的字数范围内，并避免放置公式、图表和参考文献编号。

Kresling 折纸结构具有轴向伸缩、扭转耦合和多稳态等特征。后续可将实际研究中的理论建模、有限元分析、样机实验和控制结果依次替换到正文各章中，同时保持本项目的封面、目录、章节层级与参考文献设置。

\cnkeywords{Kresling 折纸；柔性机构；力学建模；硕士学位论文}

\cleardoublepage
\chapter*{ABSTRACT}
\addcontentsline{toc}{chapter}{ABSTRACT}

This document provides a modular master's thesis layout reconstructed from the graduate-thesis formatting specification of Zhejiang University of Technology. The present abstract is placeholder text for checking typography, headers, page numbering, and keyword formatting. It should later be replaced with an English abstract consistent with the Chinese version and describing the objective, methods, principal results, and conclusions of the actual research.

\enkeywords{Kresling origami, compliant mechanism, mechanical modeling, master's thesis}
